\chapter{Perron interpolation and the geodesic equation}
\label{chap:hermitian-rays}

Positivity on a product with a strip imposes convexity in the real parameter. Taking the largest such path below two prescribed endpoint potentials gives Perron interpolation. When the endpoints have the same singularity type, elementary barriers ensure that the envelope attains them and is Lipschitz in time. The smooth product equation has the familiar homogeneous Monge--Amp\`ere form, but the geodesic equation for the $L^2$ metric contains additional derivatives of the background. Constant translations of a potential already detect this difference.

\section{Positivity on a strip}\label{sec:strip-positivity}

For a nonempty open interval $I\subset\R$, write
\[
 S_I=\{z\in\C:\Real z\in I\},\qquad
 \pi:X\times S_I\longrightarrow X.
\]
A family $u_t\in\PSH(X,\theta)$, $t\in I$, is a \emph{subgeodesic} if its complexification $U(x,z)=u_{\Real z}(x)$ is quasi-plurisubharmonic and satisfies $\pi^*\theta+\ddc U\geq0$ on the product. For real $a<b$ and $u_a,u_b\in\PSH(X,\theta)$, a subgeodesic on $(a,b)$ from $u_a$ to $u_b$ additionally converges to those endpoints in $L^1(X)$. Here $L^1$ convergence is measured against any fixed smooth positive volume form.

\begin{proposition}\label{prop:hermitian-subgeodesic-calculus}
Subgeodesics are preserved by adding an affine function of $t$, by a positive affine reparametrization of time, by finite maxima, and by convex combinations. For $h\in C^\infty(X,\mathbb R)$, replacing $\theta$ by $\theta+\ddc h$ and $u_t$ by $u_t-h$ also preserves them.
\end{proposition}
\begin{proof}
Affine functions of $\Real z$ have zero complex Hessian, and a positive affine change of time is induced by a holomorphic change of strip coordinates. Convex combinations preserve the current inequality, and smooth gauge changes preserve the current itself. Maxima are covered by \cref{lem:max-psh}.
\end{proof}

For suprema of subgeodesics we use the local envelope theorem of \cref{sec:obstacle-envelopes} on $X\times S_I$. Regularization is taken on this product; it preserves invariance in $\Im z$.

\section{Interior continuity and endpoint control}

The vertical convexity in \cref{lem:invariant-strip-convexity} bounds a path above by its chords. The minimum principle gives an integrable lower bound on every bounded interior time interval, so pointwise convexity also yields $L^1$ continuity. At an endpoint, the chord estimate converts $L^1$ convergence into a joint upper bound in space and time.

\begin{proposition}
\label{prop:subgeodesic-endpoint-control}
A subgeodesic is continuous in $L^1(X)$ on the interior of its parameter interval. If $a<b$ are finite endpoints and its $L^1$ endpoint limits are $u_a,u_b\in\PSH(X,\theta)$, then
\begin{equation}\label{eq:hermitian-chord-bound}
 u_t\leq\frac{b-t}{b-a}u_a+\frac{t-a}{b-a}u_b,
 \qquad a<t<b.
\end{equation}
Moreover the boundary upper limits hold jointly in space and time:
\[
 \limsup_{\substack{y\to x\\t\downarrow a}}u_t(y)\leq u_a(x),
 \qquad
 \limsup_{\substack{y\to x\\t\uparrow b}}u_t(y)\leq u_b(x).
\]
\end{proposition}
\begin{proof}
On a compact interior interval, choose two surrounding times. Convexity bounds the potentials above by the chord of the two corresponding potentials, hence by a constant. \Cref{prop:hermitian-kiselman} supplies one genuine integrable lower bound on the interval between the surrounding times. For almost every $x$, pointwise continuity follows from finite convexity. Dominated convergence now proves interior $L^1$ continuity.

For the endpoint inequality, choose sequences approaching $a,b$ along which the stipulated $L^1$ convergence holds almost everywhere. Apply convexity between these times and pass to the limit. For each $t$, this proves \eqref{eq:hermitian-chord-bound} almost everywhere in $x$, hence everywhere by \cref{lem:qpsh-almost-everywhere-order}.

Near $a$, the term involving $u_b$ is bounded above by a constant times its vanishing coefficient. Upper semicontinuity of $u_a$ then gives the joint upper limit. If $u_a(x)=-\infty$, first bound $u_a$ above near $x$ by an arbitrary real number and then let that number tend to $-\infty$. The argument at $b$ is the same.
\end{proof}

\Needspace{18\baselineskip}
\section{Perron segments}\label{sec:perron-segments}

The endpoint upper bounds are stable under taking the largest subgeodesic. Attaining the boundary values requires a lower barrier as well. Fix $u_0,u_1\in\PSH(X,\theta)$ and work on $[0,1]$. Call a subgeodesic $v_t$ \emph{admissible for $u_0,u_1$} if its joint boundary upper limits are at most $u_0,u_1$, respectively. Define the Perron envelope on the product by
\[
 U=\left(\sup\{V:V\text{ is an admissible complexification}\}\right)^*.
\]
Its slices are denoted $u_t^{\mathrm P}$. If it is a subgeodesic attaining the prescribed endpoints in $L^1$, we call it the \emph{Perron segment} between them and set $u_0^{\mathrm P}=u_0$, $u_1^{\mathrm P}=u_1$. We use the same definition on any finite interval by affine reparametrization.

\begin{proposition}\label{prop:hermitian-perron}
If there is an admissible competitor, the Perron envelope is genuine, is itself admissible, and is the largest admissible subgeodesic. If at least one competitor attains $u_0,u_1$ in $L^1$, the envelope attains them as well. In that case its restriction to every subinterval is the Perron envelope of its own endpoints.
\end{proposition}
\begin{proof}
Vertical convexity and the boundary inequalities bound every competitor by
$(1-t)u_0+t u_1$. This is jointly upper semicontinuous on the interior and locally bounded above, so the same bound holds after regularization of the supremum. The local envelope theorem gives a product potential. It dominates a genuine competitor, so every slice is genuine. The chord bound also gives the joint boundary inequalities, as in the proof of \cref{prop:subgeodesic-endpoint-control}. Thus the envelope belongs to the family defining it, and is its largest member.

If a competitor $w_t$ has the required $L^1$ endpoints, then
\[
 w_t\leq u_t^{\mathrm P}\leq(1-t)u_0+t u_1.
\]
Both outer functions converge in $L^1$ to $u_0$ as $t\downarrow0$. Almost everywhere,
\[
 |u_t^{\mathrm P}-u_0|\leq |w_t-u_0|+t|u_1-u_0|,
\]
which proves the same convergence for the envelope. The argument at $1$ is identical.

For the restriction assertion, fix $0\leq a<b\leq1$ and let $v_t$ be the Perron envelope between $u_a^{\mathrm P}$ and $u_b^{\mathrm P}$. The restricted original path is a competitor attaining these endpoints, by \cref{prop:subgeodesic-endpoint-control}. Hence $v_t$ attains them and $v_t\geq u_t^{\mathrm P}$ on $(a,b)$. By \cref{lem:qpsh-gluing}, its complexification can be pasted to $U$ outside the smaller strip: the boundary upper limits are at most the corresponding slices of $U$, and the maximum inside equals the new envelope. The pasted path is admissible for $u_0,u_1$. Maximality of $u_t^{\mathrm P}$ forces equality with $v_t$ on $(a,b)$.
\end{proof}

\Needspace{15\baselineskip}
A bounded difference between the endpoints supplies the required lower barrier explicitly.

\begin{theorem}
\label{thm:ordinary-type-perron}
Suppose $u_0\sim u_1$ and choose $C\geq0$ with
$u_0-C\leq u_1\leq u_0+C$. The Perron segment exists and attains both endpoints. Every slice has their ordinary singularity type, and
\[
 u_0-Ct\leq u_t^{\mathrm P}\leq u_0+Ct,
 \qquad
 u_1-C(1-t)\leq u_t^{\mathrm P}\leq u_1+C(1-t).
\]
Away from the common pole set,
\[
 |u_t^{\mathrm P}-u_s^{\mathrm P}|\leq C|t-s|,
 \qquad0\leq s,t\leq1.
\]
\end{theorem}
\begin{proof}
The function
\[
 w_t=\max\{u_0-Ct,\ u_1-C(1-t)\}
\]
is a subgeodesic by \cref{prop:hermitian-subgeodesic-calculus}. The bounded difference gives $w_0=u_0,w_1=u_1$. Upper semicontinuity gives its joint boundary upper limits, and the same bounded difference gives $|w_t-u_0|\leq Ct$ and $|w_t-u_1|\leq C(1-t)$ off the common poles. Thus $w_t$ attains both endpoints in $L^1$, and \cref{prop:hermitian-perron} applies.

The lower bounds follow from this competitor and the upper bounds from the chord estimate. At a point off the common poles, the lower bound at $t=0$ gives a right derivative of at least $-C$, while the lower bound at $t=1$ gives a left derivative of at most $C$. Every intermediate secant slope of a finite convex function lies between these two endpoint slopes. Hence all secant slopes belong to $[-C,C]$.
\end{proof}

For endpoints of different types, a common lower potential gives an admissible constant path. To obtain a segment attaining those endpoints, the lower barrier must also approach the boundary values, as in \cref{prop:hermitian-perron}.

\Needspace{25\baselineskip}
\section{The smooth product equation}

For a smooth path with positive spatial form, the product positivity condition reduces to a scalar inequality by taking a Schur complement. Let $u_t$ be smooth with $\eta_t=\theta+\ddc u_t>0$, and let dots denote time derivatives. In local holomorphic coordinates $x_1,\ldots,x_n$, write
\[
 \theta=\frac{\mathrm i}{2\pi}\sum_{i,j=1}^n
        g_{i\bar j}\,\mathrm dx_i\wedge\mathrm d\bar x_j,
 \qquad H=(g_{i\bar j}+\partial_i\partial_{\bar j}u),
 \qquad \xi=(\partial_i\dot u)_{i=1}^n.
\]
Here $\xi$ is a column vector and $\xi^*$ is its conjugate transpose. Since $t=\Real z$, the product form $\pi^*\theta+\ddc U$ is represented by
\[
 \begin{pmatrix}
 H&\tfrac12\xi\\
 \tfrac12\xi^*&\tfrac14\ddot u
 \end{pmatrix}.
\]

\Needspace{12\baselineskip}
\begin{proposition}\label{prop:hermitian-schur-geodesic}
The product form is semipositive exactly when
\[
 \ddot u\geq\xi^*H^{-1}\xi.
\]
Its $(n+1)$st power vanishes exactly when equality holds. Equivalently, the homogeneous product equation is
\begin{equation}\label{eq:hermitian-hcma-smooth}
 \ddot u\,\eta_t^n
 =n\,\mathrm{d}\dot u\wedge\dc\dot u\wedge\eta_t^{n-1}.
\end{equation}
\end{proposition}
\begin{proof}
The Schur complement of $H$ is $(\ddot u-\xi^*H^{-1}\xi)/4$. Since $H>0$, the block matrix is semipositive exactly when this complement is nonnegative, and its determinant vanishes exactly when the complement vanishes.

For a smooth real function $f$, put $\xi_f=(\partial_i f)_{i=1}^n$. Then
\[
 \mathrm{d}f\wedge\dc f=\frac{\mathrm i}{2\pi}\partial f\wedge\bar\partial f,
 \qquad
 n\frac{\mathrm{d}f\wedge\dc f\wedge\eta_t^{n-1}}{\eta_t^n}
 =\operatorname{tr}(H^{-1}\xi_f\xi_f^*)=\xi_f^*H^{-1}\xi_f.
\]
Apply this identity with $f=\dot u$ to obtain \eqref{eq:hermitian-hcma-smooth}.
\end{proof}

The equation uses only the values of the spatial form $\eta_t$, whereas the variation of an integral involving $\eta_t^n$ also involves its derivatives. This distinction produces the torsion terms below. Establishing the homogeneous equation for nonsmooth Perron segments would require a separate weak-solution argument.

\Needspace{25\baselineskip}
\section{The \texorpdfstring{$L^2$}{L2} action and torsion}

Fix a Hermitian form $\omega$. On $\HH_\omega$, the Monge--Amp\`ere volume of each potential defines the $L^2$ metric
\[
 G_u(v,w)=\int_Xvw\,\omega_u^n,
 \qquad v,w\in C^\infty(X,\mathbb R).
\]
For a smooth path $u_t\in\HH_\omega$, $0\leq t\leq1$, its action is
\[
 \mathscr S(u)=\frac12\int_0^1\int_X\dot u_t^{\,2}\,\omega_{u_t}^n\,\mathrm{d}t.
\]
Varying the velocity gives the acceleration term; varying the volume form and integrating by parts gives the spatial gradient and torsion terms. Stationarity under smooth variations fixing both endpoints is the formal geodesic equation for this metric.

\Needspace{10\baselineskip}
\begin{theorem}\label{thm:hermitian-ltwo-stationarity}
Put $\eta_t=\omega_{u_t}$ and $T_t=\eta_t^{n-1}$. The path is stationary for $\mathscr S$ if and only if
\begin{equation}\label{eq:hermitian-ltwo-geodesic}
 \begin{split}
 \ddot u_t\,\eta_t^n={}&n\,\mathrm{d}\dot u_t\wedge\dc\dot u_t\wedge T_t\\
 &+n\dot u_t\bigl(\mathrm{d}\dot u_t\wedge\dc T_t
                  -\dc\dot u_t\wedge \mathrm{d}T_t\bigr)
   +\frac n2\dot u_t^{\,2}\ddc T_t.
 \end{split}
\end{equation}
All differentials on the right act only in the $X$ variable.
\end{theorem}
\begin{proof}
Let $v_t$ be a smooth variation with $v_0=v_1=0$, and put $f=\dot u_t$. Differentiating the action in its variation parameter gives
\[
 \delta\mathscr S
 =\int_0^1\left(\int_X f\dot v_t\,\eta_t^n
       +\frac n2\int_X f^2\ddc v_t\wedge T_t\right)\mathrm{d}t.
\]
Integration in time uses
$\partial_t(\eta_t^n)=n\ddc f\wedge T_t$. Integration on $X$ uses \cref{lem:ibp} with the entire form $f^2T_t$. Thus
\[
 \delta\mathscr S
 =\int_0^1\int_Xv_t\left(-\ddot u_t\eta_t^n
       -nf\ddc f\wedge T_t+\frac n2\ddc(f^2T_t)\right)\mathrm{d}t.
\]
The full product rule is
\[
 \begin{split}
 \ddc(f^2T_t)={}&2f\ddc f\wedge T_t
       +2\mathrm{d}f\wedge\dc f\wedge T_t\\
 &+2f\bigl(\mathrm{d}f\wedge\dc T_t-\dc f\wedge \mathrm{d}T_t\bigr)
       +f^2\ddc T_t.
 \end{split}
\]
The first term cancels $-nf\ddc f\wedge T_t$. The remaining expression is the right side of \eqref{eq:hermitian-ltwo-geodesic} minus its left side, integrated against $v_t$. Arbitrary compactly supported variations in time and space and the fundamental lemma of the calculus of variations prove necessity. Substitution proves sufficiency. Small variations are admissible because the original path is strictly positive on the compact parameter interval.
\end{proof}

For a K\"ahler background, $\mathrm{d}T_t=\dc T_t=0$ and the equation reduces to \eqref{eq:hermitian-hcma-smooth}. The Gauduchon condition on $\eta_t$ removes the last term but leaves the terms involving $\mathrm{d}T_t$ and $\dc T_t$. Constant translations isolate the last term.

\Needspace{10\baselineskip}
\begin{corollary}
\label{cor:translations-gauduchon}
Fix $u\in\HH_\omega$ and a nonzero constant $a$. The path $u_t=u+at$ is a Perron segment and solves the smooth homogeneous product equation. It is stationary for the $L^2$ action if and only if
\[
 \ddc(\omega_u^{n-1})=0.
\]
For $n\geq2$, this says precisely that $\omega_u$ is Gauduchon.
\end{corollary}
\begin{proof}
The path is admissible and equals the chord between its endpoints, so no admissible competitor can be larger. It is therefore the Perron segment. Its time acceleration and the spatial derivatives of its velocity vanish, which proves the homogeneous product equation. In \eqref{eq:hermitian-ltwo-geodesic}, the sole remaining term is
$\frac n2a^2\ddc(\omega_u^{n-1})$. This vanishes exactly under the displayed condition.
\end{proof}

For the Gauduchon background and its perturbation in \cref{ex:gauduchon-persistence-torus}, constant translation through the perturbation is consequently a smooth Perron segment which is not stationary for the $L^2$ action. On a surface the Gauduchon condition persists under potential perturbations, as explained in \cref{prop:gauduchon-potential-persistence}, but it still supplies no cancellation of the first-derivative terms in \eqref{eq:hermitian-ltwo-geodesic}. The Legendre-transform constructions of the next chapter use Perron maximality; their relation to a complete metric on Hermitian singularity classes remains a separate problem.
